\chapter{TỔNG QUAN LÝ THUYẾT VÀ CƠ SỞ NGHIÊN CỨU}
\section{ Tổng quan về ô nhiễm không khí và $PM_{2.5}$}
\subsection{Ô nhiễm không khí tại Việt Nam}
Ô nhiễm không khí đang là một trong những thách thức môi trường nghiêm trọng nhất trên phạm vi toàn cầu. Theo ước tính của Tổ chức Y tế Thế giới (WHO), ô nhiễm không khí gây ra khoảng 7 triệu ca tử vong sớm mỗi năm, trong đó các thành phố lớn tại châu Á có nồng độ $PM_{2.5}$ vượt gấp 10 lần mức khuyến nghị \cite{krzyzanowski2014megacities}. Tại Việt Nam, ô nhiễm không khí được xếp vào nhóm \textbf{5 yếu tố hàng đầu} gây gánh nặng bệnh tật và tử vong sớm trong năm 2019, chỉ đứng sau huyết áp cao, hút thuốc, tiểu đường và rủi ro dinh dưỡng \cite{hoang2023health}.

Quá trình đô thị hóa nhanh chóng tại Việt Nam là một trong những nguyên nhân chính dẫn đến sự gia tăng ô nhiễm không khí. Nghiên cứu hồi quy trên dữ liệu 63 tỉnh thành cho thấy mỗi 1\% gia tăng tỷ lệ đô thị hóa tương ứng với mức tăng \textbf{0,54 $\mu g/m^3$} nồng độ $PM_{2.5}$ \cite{bui2025urbanization}. Mật độ dân cư cao tại các đô thị kéo theo sự gia tăng phương tiện giao thông, hoạt động xây dựng và sản xuất công nghiệp, trong khi hạ tầng quản lý và xử lý ô nhiễm chưa phát triển tương xứng. Báo cáo của RTI International (2024) chỉ ra rằng khoảng 70\% nước thải công nghiệp tại Việt Nam chưa được xử lý đúng quy trình, phản ánh một bức tranh tổng thể về áp lực ô nhiễm lên môi trường \cite{rti2024pollution}.

Xét về nguồn phát thải, giao thông đóng vai trò chủ đạo trong việc phát sinh các chất ô nhiễm như $NO_2$, $CO$ và $PM_{2.5}$ tại Hà Nội. Nghiên cứu của Ho et al. (2025) sử dụng mô hình CMAQ kết hợp dữ liệu phát thải giao thông cho thấy nồng độ các chất ô nhiễm từ nguồn này tăng cao đáng kể vào mùa đông \cite{ho2025assessment}. Bên cạnh giao thông, các hoạt động công nghiệp và làng nghề truyền thống (gốm, hương, rèn) cũng đóng góp một lượng đáng kể chất ô nhiễm vào khí quyển \cite{hoang2023health}. Đặc biệt, hoạt động đốt rơm rạ và phụ phẩm nông nghiệp làm gia tăng nồng độ $NO_2$ tại các vùng như Duyên hải Bắc Trung Bộ, Đông Nam Bộ và Đồng bằng sông Cửu Long \cite{hoang2023health}. Ngoài các nguồn phát thải sơ cấp, $PM_{2.5}$ còn được hình thành thông qua các phản ứng hóa học thứ cấp trong khí quyển, khi $SO_2$, $NO_2$ và các hợp chất hữu cơ dễ bay hơi (VOCs) chuyển hóa thành aerosol sulfate và nitrate \cite{le2024environmental}.

Phân bố ô nhiễm không khí tại Việt Nam có sự khác biệt rõ rệt giữa các vùng miền. Khu vực \textbf{Đồng bằng sông Hồng} ghi nhận nồng độ $PM_{2.5}$ và $NO_2$ cao nhất cả nước, với nồng độ $PM_{2.5}$ trung bình năm dao động từ 12,8 đến 40,8 $\mu g/m^3$ trong năm 2019 \cite{hoang2023health}. Riêng Hà Nội, ước tính có khoảng 315 ca tử vong có thể tránh được mỗi năm nếu giảm nồng độ $NO_2$ xuống mức thời kỳ COVID-19, trong bối cảnh mật độ dân số khu vực này đạt khoảng 1.000 người/$km^2$ \cite{hoang2023health}. Ngược lại, khu vực \textbf{miền Nam} ghi nhận mức ô nhiễm thấp hơn đáng kể, với số ca tử vong tránh được từ giảm $NO_2$ chỉ dao động từ 3 đến 37 ca tùy tỉnh \cite{hoang2023health}. Sự chênh lệch này được giải thích bởi đặc điểm khí hậu: mùa đông miền Bắc thường xuất hiện hiện tượng nghịch nhiệt, tầng trộn thấp và gió yếu, khiến các chất ô nhiễm bị giữ lại gần mặt đất và tích tụ \cite{dam2024pm25, ho2025assessment}.

Biến động theo mùa là một đặc trưng quan trọng của ô nhiễm không khí tại Việt Nam. Vào giai đoạn \textbf{thu--đông}, nồng độ $PM_{2.5}$ tăng cao do hiện tượng nghịch nhiệt làm giảm chiều cao tầng trộn, kết hợp với việc tăng sử dụng nhiên liệu hóa thạch cho sưởi ấm \cite{dam2024pm25, ho2025assessment}. Mô phỏng từ mô hình CMAQ kết hợp WRF cho thấy nồng độ $PM_{2.5}$ tại Hà Nội ở độ cao 30 mét đạt 50--84 $\mu g/m^3$, vượt đáng kể Quy chuẩn kỹ thuật quốc gia QCVN 05:2013 (50 $\mu g/m^3$) và khuyến nghị của WHO (25 $\mu g/m^3$) \cite{nguyen2024cmaq}. Nồng độ này giảm dần theo độ cao, chỉ còn 34--40 $\mu g/m^3$ ở độ cao 320 mét, cho thấy ô nhiễm tập trung mạnh ở tầng thấp nơi con người sinh sống \cite{dam2024pm25}. Vào \textbf{mùa hè}, nhờ lượng mưa lớn, tầng trộn cao và điều kiện khuếch tán thuận lợi, nồng độ ô nhiễm giảm đáng kể \cite{dam2024pm25}.

Giai đoạn giãn cách xã hội do đại dịch COVID-19 đã cung cấp một ``thí nghiệm tự nhiên'' quý giá, minh chứng rõ ràng cho mối quan hệ nhân quả giữa hoạt động con người và chất lượng không khí. Trong thời gian phong tỏa tháng 4 năm 2020 tại Hà Nội, nồng độ $PM_{2.5}$ \textbf{giảm 75,8\%} và $NO_2$ \textbf{giảm 41,8\%} so với cùng kỳ \cite{hoang2023health}. Trên phạm vi toàn quốc 63 tỉnh thành, mức giảm $PM_{2.5}$ trung bình năm dao động từ 5,7\% đến 10,3\% \cite{hoang2023health}. Nếu duy trì được mức chất lượng không khí tương đương giai đoạn phong tỏa, ước tính có thể tránh được \textbf{3.807 ca tử vong} do $PM_{2.5}$ (tiết kiệm khoảng 793 triệu USD) và \textbf{2.451 ca tử vong} do $NO_2$ (tiết kiệm khoảng 510,6 triệu USD) \cite{hoang2023health}. Những con số này nhấn mạnh tầm quan trọng của việc phát triển các hệ thống dự báo chất lượng không khí chính xác, nhằm hỗ trợ cảnh báo sớm và cung cấp cơ sở khoa học cho các chính sách kiểm soát ô nhiễm tại Việt Nam.

\subsection{Bụi mịn $PM_{2.5}$}

Bụi mịn $PM_{2.5}$ (Particulate Matter 2.5) là thuật ngữ chỉ các hạt vật chất lơ lửng trong không khí có đường kính khí động học \textbf{nhỏ hơn hoặc bằng 2,5 $\mu m$}. Do kích thước cực nhỏ, $PM_{2.5}$ có khả năng xâm nhập sâu vào phế nang phổi, thậm chí đi vào máu và các cơ quan nội tạng, gây ra nhiều bệnh lý nghiêm trọng về hô hấp, tim mạch và ung thư \cite{le2024environmental}. Tổ chức Y tế Thế giới (WHO) khuyến nghị nồng độ $PM_{2.5}$ trung bình năm không nên vượt quá 5 $\mu g/m^3$ (hướng dẫn năm 2021) \cite{WHO_2021_AQG}, tuy nhiên hầu hết các đô thị lớn tại châu Á, bao gồm Việt Nam, đều vượt xa ngưỡng này \cite{krzyzanowski2014megacities}.

Về nguồn gốc hình thành, $PM_{2.5}$ được phân thành hai loại chính. \textbf{$PM_{2.5}$ sơ cấp} là các hạt bụi được phát thải trực tiếp vào khí quyển từ quá trình đốt cháy nhiên liệu hóa thạch (xăng, dầu diesel, than đá), khí thải công nghiệp, đốt sinh khối (rơm rạ, gỗ), bụi đường và hoạt động xây dựng \cite{ho2025assessment, bui2025urbanization}. \textbf{$PM_{2.5}$ thứ cấp} được hình thành thông qua các phản ứng hóa học trong khí quyển, khi các chất tiền thân dạng khí như $SO_2$, $NO_x$ ($NO$ và $NO_2$), amoniac ($NH_3$) và các hợp chất hữu cơ dễ bay hơi (VOCs) trải qua quá trình oxy hóa, ngưng tụ và chuyển đổi pha, tạo thành các hạt aerosol sulfate ($SO_4^{2-}$), nitrate ($NO_3^{-}$) và carbon hữu cơ thứ cấp (SOA) \cite{le2024environmental}. Tỷ lệ đóng góp giữa $PM_{2.5}$ sơ cấp và thứ cấp phụ thuộc vào đặc điểm nguồn phát thải và điều kiện khí tượng của từng khu vực, trong đó tại các đô thị lớn như Hà Nội, $PM_{2.5}$ thứ cấp có thể chiếm tỷ trọng đáng kể do mật độ phát thải $NO_x$ và $SO_2$ từ giao thông và công nghiệp cao \cite{ho2025assessment}.

Nồng độ $PM_{2.5}$ trong khí quyển chịu ảnh hưởng mạnh mẽ bởi các yếu tố khí tượng. \textbf{Nhiệt độ} ảnh hưởng đến cấu trúc nhiệt của khí quyển; khi nhiệt độ bề mặt giảm nhanh hơn các tầng trên, hiện tượng nghịch nhiệt (temperature inversion) hình thành, tạo ra một ``lớp nắp'' ngăn cản sự khuếch tán thẳng đứng của chất ô nhiễm \cite{dam2024pm25}. Nghiên cứu tại Hà Nội cho thấy vào mùa đông, hiện tượng nghịch nhiệt làm giảm đáng kể chiều cao tầng trộn, khiến nồng độ $PM_{2.5}$ ở tầng thấp (30 mét) đạt 50--84 $\mu g/m^3$, trong khi ở độ cao 320 mét chỉ còn 34--40 $\mu g/m^3$ \cite{dam2024pm25, nguyen2024cmaq}. \textbf{Độ ẩm tương đối} đóng vai trò kép: độ ẩm cao thúc đẩy phản ứng hóa học tạo $PM_{2.5}$ thứ cấp (đặc biệt quá trình oxy hóa $SO_2$ thành sulfate trong pha nước), nhưng khi đạt đến mức bão hòa, mưa lại đóng vai trò rửa trôi (washout) bụi mịn khỏi khí quyển \cite{dam2024pm25}. \textbf{Tốc độ và hướng gió} quyết định khả năng vận chuyển và phân tán các chất ô nhiễm; gió yếu khiến $PM_{2.5}$ tích tụ tại nguồn, trong khi gió mạnh giúp pha loãng nhưng cũng có thể vận chuyển ô nhiễm xuyên biên giới từ các khu vực công nghiệp lân cận \cite{ho2025assessment}.

Biến động theo mùa của $PM_{2.5}$ tại Việt Nam thể hiện đặc trưng rõ rệt. Vào \textbf{mùa đông} (tháng 11 đến tháng 3), nồng độ $PM_{2.5}$ đạt mức cao nhất trong năm do sự kết hợp của ba yếu tố: (i) hiện tượng nghịch nhiệt thường xuyên, (ii) gió mùa Đông Bắc vận chuyển ô nhiễm từ các khu công nghiệp phía Bắc, và (iii) gia tăng phát thải từ hoạt động sưởi ấm bằng nhiên liệu hóa thạch \cite{dam2024pm25, ho2025assessment}. Ngược lại, vào \textbf{mùa hè} (tháng 5 đến tháng 9), lượng mưa lớn kết hợp với tầng trộn cao và bức xạ mặt trời mạnh tạo điều kiện khuếch tán thuận lợi, giúp nồng độ $PM_{2.5}$ giảm đáng kể \cite{dam2024pm25}. Phân tích của Ho et al. (2025) trên dữ liệu mô phỏng CMAQ tại Hà Nội cho thấy nồng độ các chất ô nhiễm từ nguồn giao thông cũng tuân theo biến động mùa tương tự, với đỉnh vào mùa đông và đáy vào mùa hè \cite{ho2025assessment}. Đặc điểm biến động mùa vụ này đặt ra yêu cầu quan trọng đối với các mô hình dự báo: hệ thống phải có khả năng nắm bắt tính mùa vụ (seasonality) và thích ứng với các điều kiện khí tượng thay đổi để đảm bảo độ chính xác dự báo trong mọi thời điểm của năm.

\subsection{Tác động của $PM_{2.5}$ đến sức khỏe con người}

Bụi mịn $PM_{2.5}$ được xem là chất ô nhiễm không khí nguy hiểm nhất đối với sức khỏe con người, do kích thước cực nhỏ cho phép các hạt bụi vượt qua hàng rào bảo vệ tự nhiên của hệ hô hấp. Trong khi các hạt bụi thô ($PM_{10}$) bị giữ lại ở mũi và đường thở trên, $PM_{2.5}$ có khả năng xâm nhập sâu vào tận phế nang phổi, từ đó khuếch tán qua màng phế nang -- mao mạch vào hệ tuần hoàn máu và phân bố đến các cơ quan nội tạng như tim, gan, thận và não \cite{le2024environmental}. 

Về tác động lên \textbf{hệ hô hấp}, tiếp xúc lâu dài với $PM_{2.5}$ làm tăng nguy cơ mắc các bệnh viêm phế quản mãn tính, hen suyễn, bệnh phổi tắc nghẽn mãn tính (COPD) và ung thư phổi. Các hạt bụi mịn khi lắng đọng trong phế nang kích hoạt phản ứng viêm, tạo ra các gốc tự do (ROS -- Reactive Oxygen Species) gây tổn thương mô phổi và suy giảm chức năng hô hấp \cite{le2024environmental}. Đối với \textbf{hệ tim mạch}, nghiên cứu tổng quan của Le et al. (2024) chỉ ra rằng $PM_{2.5}$ và $NO_2$ có mối liên hệ trực tiếp với các bệnh lý tim mạch bao gồm nhồi máu cơ tim, đột quỵ và xơ vữa động mạch. Cơ chế bệnh sinh bao gồm phản ứng viêm hệ thống, stress oxy hóa và rối loạn chức năng nội mô mạch máu \cite{le2024environmental}. Ngoài ra, các nghiên cứu gần đây còn ghi nhận mối liên hệ giữa phơi nhiễm $PM_{2.5}$ với các bệnh lý thần kinh, tiểu đường type 2 và suy giảm chức năng nhận thức \cite{krzyzanowski2014megacities}.

Trên phạm vi toàn cầu, ô nhiễm không khí gây ra khoảng \textbf{7 triệu ca tử vong sớm} mỗi năm, trong đó $PM_{2.5}$ chiếm phần lớn gánh nặng bệnh tật \cite{krzyzanowski2014megacities}. Tại Việt Nam, ô nhiễm không khí được xếp trong nhóm 5 yếu tố nguy cơ hàng đầu gây gánh nặng bệnh tật năm 2019, ngang hàng với các yếu tố truyền thống như huyết áp cao và hút thuốc \cite{hoang2023health}. Nghiên cứu định lượng của Hoang et al. (2023) trên dữ liệu 63 tỉnh thành đã ước tính cụ thể thiệt hại sức khỏe từ ô nhiễm không khí: nếu nồng độ $PM_{2.5}$ được duy trì ở mức thời kỳ giãn cách COVID-19, Việt Nam có thể tránh được \textbf{3.807 ca tử vong} mỗi năm, tương đương giá trị kinh tế khoảng \textbf{793 triệu USD}. Tương tự, việc giảm $NO_2$ có thể tránh được \textbf{2.451 ca tử vong}, tiết kiệm khoảng \textbf{510,6 triệu USD} \cite{hoang2023health}. 

Phân bố thiệt hại sức khỏe có sự chênh lệch rõ rệt giữa các vùng miền, phản ánh sự khác biệt về mức độ ô nhiễm. Khu vực Đồng bằng sông Hồng, đặc biệt là Hà Nội với mật độ dân số khoảng 1.000 người/$km^2$, chịu tác động nặng nề nhất với ước tính \textbf{315 ca tử vong} có thể tránh được mỗi năm chỉ riêng từ việc giảm $NO_2$ \cite{hoang2023health}. Trong khi đó, các tỉnh miền Nam với nồng độ ô nhiễm thấp hơn ghi nhận con số chỉ từ 3 đến 37 ca tùy địa phương \cite{hoang2023health}. Sự phân bố không đồng đều này cho thấy các hệ thống dự báo chất lượng không khí cần có khả năng \textbf{phân giải theo không gian} (spatial resolution), đặc biệt tại các khu vực đô thị mật độ cao, nhằm phục vụ hiệu quả cho công tác cảnh báo sức khỏe cộng đồng.

Những bằng chứng trên khẳng định rằng ô nhiễm $PM_{2.5}$ không chỉ là vấn đề môi trường mà còn là \textbf{vấn đề y tế công cộng và kinh tế xã hội} có quy mô lớn. Điều này đặt ra nhu cầu cấp thiết về việc xây dựng các hệ thống dự báo nồng độ $PM_{2.5}$ chính xác, có khả năng cung cấp cảnh báo sớm trước 1 đến 24 giờ, giúp cơ quan quản lý ban hành khuyến cáo sức khỏe kịp thời và người dân chủ động bảo vệ bản thân, đặc biệt trong các đợt ô nhiễm nghiêm trọng vào mùa đông.

\section{Cơ sở lý thuyết kỹ thuật}
\subsection{Bài toán dự báo chuỗi thời gian}

\textbf{Dự báo chuỗi thời gian} (time series forecasting) là bài toán ước lượng giá trị tương lai của một biến số dựa trên chuỗi các quan sát được ghi nhận theo trình tự thời gian. Cho trước chuỗi thời gian $\mathbf{x} = \{x_1, x_2, \ldots, x_T\}$ với $T$ bước thời gian đã quan sát, mục tiêu của bài toán là xây dựng hàm ánh xạ $f$ sao cho:
\begin{equation}
    \hat{x}_{T+h} = f(x_{T-L+1}, x_{T-L+2}, \ldots, x_T; \boldsymbol{\theta})
\end{equation}
trong đó $L$ là độ dài cửa sổ quan sát (lookback window), $h$ là tầm nhìn dự báo (forecast horizon) và $\boldsymbol{\theta}$ là tập tham số của mô hình. Trong bài toán đa biến (multivariate), đầu vào mở rộng thành ma trận $\mathbf{X} \in \mathbb{R}^{L \times d}$ với $d$ đặc trưng (features), bao gồm cả biến mục tiêu và các biến ngoại sinh (exogenous variables) như nhiệt độ, độ ẩm, tốc độ gió \cite{Chen_2026, liu2024itransformer}.

\subsubsection{Dự báo đơn bước và đa tầm nhìn}

Dựa trên số bước thời gian cần dự báo, bài toán được phân thành hai dạng chính. \textbf{Dự báo đơn bước} (single-step forecasting) chỉ ước lượng giá trị tại một thời điểm duy nhất $T+1$, phù hợp cho các ứng dụng cần phản ứng nhanh như cảnh báo ô nhiễm tức thì. Nghiên cứu của Silveira et al. (2025) cho thấy LSTM và GRU đạt hiệu suất cao trên bài toán 1-step, tuy nhiên sai số tăng nhanh khi mở rộng sang 20-step \cite{silveira2025lstmgru}. \textbf{Dự báo đa tầm nhìn} (multi-horizon forecasting) ước lượng đồng thời nhiều bước tương lai $\{\hat{x}_{T+1}, \hat{x}_{T+3}, \ldots, \hat{x}_{T+H}\}$, phục vụ lập kế hoạch trung và dài hạn \cite{armstrong1992mape}. Bài toán multi-horizon đặt ra thách thức lớn hơn do sai số tích lũy (error accumulation) và sự suy giảm tương quan giữa đầu vào và đầu ra khi tầm nhìn tăng. Trong đồ án này, hệ thống được thiết kế để benchmark trên \textbf{5 tầm nhìn}: $T+1$, $T+3$, $T+6$, $T+12$ và $T+24$ giờ, tương ứng với các kịch bản từ cảnh báo tức thì đến dự báo trong ngày.

\subsubsection{Đặc điểm chuỗi thời gian $PM_{2.5}$}

Chuỗi thời gian nồng độ $PM_{2.5}$ sở hữu nhiều đặc trưng phức tạp, đặt ra yêu cầu cao đối với các mô hình dự báo.

\textbf{Xu hướng dài hạn (Trend).} Nồng độ $PM_{2.5}$ thể hiện xu hướng biến đổi theo thời gian, phản ánh các thay đổi cấu trúc trong phát thải và chính sách môi trường. Nghiên cứu trên dữ liệu 63 tỉnh thành Việt Nam cho thấy nồng độ $PM_{2.5}$ giảm 5,7--10,3\% trong giai đoạn 2020--2021 so với 2019 do tác động của giãn cách xã hội \cite{hoang2023health}. Tuy nhiên, quá trình đô thị hóa liên tục (mỗi 1\% tăng đô thị hóa $\rightarrow$ +0,54 $\mu g/m^3$) tạo ra xu hướng tăng dài hạn đối lập \cite{bui2025urbanization}. Sự đan xen giữa các xu hướng ngược chiều này đòi hỏi mô hình phải có khả năng phân tách và nắm bắt thành phần trend một cách linh hoạt.

\textbf{Tính mùa vụ (Seasonality).} Chuỗi $PM_{2.5}$ tại Việt Nam thể hiện tính mùa vụ rõ rệt ở nhiều chu kỳ lồng nhau. Ở \textit{chu kỳ năm}, nồng độ đạt đỉnh vào mùa đông (tháng 11--3) do nghịch nhiệt và gió mùa Đông Bắc, và đạt đáy vào mùa hè (tháng 5--9) nhờ mưa rửa trôi \cite{dam2024pm25, ho2025assessment}. Ở \textit{chu kỳ ngày}, nồng độ thường tăng cao vào giờ cao điểm giao thông (6--9h sáng và 17--20h tối) do phát thải từ phương tiện cơ giới \cite{ho2025assessment}. Ở \textit{chu kỳ tuần}, nồng độ cuối tuần thường thấp hơn ngày thường do giảm hoạt động sản xuất và giao thông. Các mô hình dự báo cần nhận diện và mã hóa được các chu kỳ mùa vụ đa tầng này.

\textbf{Tính không dừng (Non-stationarity).} Chuỗi $PM_{2.5}$ thường không thỏa mãn điều kiện dừng (stationarity), nghĩa là trung bình, phương sai và cấu trúc tự tương quan thay đổi theo thời gian. Nguồn gốc của tính không dừng bao gồm: biến đổi xu hướng do chính sách môi trường, đột biến do các sự kiện cực đoan (cháy rừng, đốt rơm rạ), và sự thay đổi điều kiện khí tượng theo mùa \cite{dam2024pm25}. Mô phỏng CMAQ tại Hà Nội cho thấy nồng độ $PM_{2.5}$ dao động trong khoảng rất rộng từ 34 đến 84 $\mu g/m^3$ chỉ trong cùng một mùa đông, phản ánh sự biến động cao \cite{nguyen2024cmaq}. Tính không dừng đặt ra thách thức lớn cho các phương pháp thống kê truyền thống (ARIMA, SARIMA) vốn yêu cầu chuỗi dừng, và là một trong những động lực thúc đẩy việc áp dụng các mô hình học sâu có khả năng học biểu diễn phi tuyến \cite{silveira2025lstmgru}.

\textbf{Phụ thuộc không gian (Spatial dependency).} Khác với nhiều bài toán chuỗi thời gian cổ điển, nồng độ $PM_{2.5}$ tại một trạm quan trắc không chỉ phụ thuộc vào lịch sử của chính nó mà còn chịu ảnh hưởng mạnh bởi các trạm lân cận thông qua quá trình vận chuyển khí quyển. Nghiên cứu của Dam et al. (2024) chỉ ra rằng nồng độ $PM_{2.5}$ có gradient rõ rệt theo chiều thẳng đứng (50--84 $\mu g/m^3$ ở 30m so với 34--40 $\mu g/m^3$ ở 320m), cho thấy ô nhiễm phân bố không đồng nhất trong không gian \cite{dam2024pm25}. Trên mặt phẳng ngang, hướng gió và vận tốc gió quyết định mô hình lan truyền ô nhiễm giữa các khu vực. Panja et al. (2025) đề xuất mô hình E-STGCN sử dụng ma trận kề (adjacency matrix) dựa trên hàm Gaussian kernel của khoảng cách Haversine để biểu diễn mối tương quan không gian giữa các trạm \cite{panja2025estgcn}. Tuy nhiên, cách tiếp cận khoảng cách tĩnh này chưa phản ánh đầy đủ tính động của quá trình vận chuyển ô nhiễm, vốn thay đổi liên tục theo điều kiện gió — đây chính là khoảng trống mà đồ án này hướng đến bằng việc đề xuất \textbf{Dynamic Graph Builder} tích hợp thông tin hướng gió và vận tốc gió vào quá trình xây dựng đồ thị không gian.

Tóm lại, sự kết hợp đồng thời của xu hướng dài hạn, tính mùa vụ đa tầng, tính không dừng và phụ thuộc không gian khiến bài toán dự báo $PM_{2.5}$ trở thành một trong những bài toán phức tạp nhất trong lĩnh vực dự báo chuỗi thời gian. Điều này đòi hỏi sự kết hợp linh hoạt giữa nhiều lớp mô hình — từ Machine Learning truyền thống đến Deep Learning và Spatio-temporal Graph Networks — nhằm khai thác hiệu quả cả thông tin thời gian lẫn không gian.

\subsection{Nhóm mô hình Machine Learning}
\subsubsection{Gradient Boosting và XGBoost}

Gradient Boosting là phương pháp học ensemble dựa trên nguyên lý \textbf{boosting}: thay vì huấn luyện một mô hình mạnh duy nhất, phương pháp này kết hợp tuần tự nhiều mô hình yếu (weak learners), trong đó mỗi mô hình mới được huấn luyện để sửa chữa sai số còn lại (residual) của tập mô hình hiện tại. Cụ thể, tại bước lặp thứ $t$, mô hình $f_t$ được thêm vào để tối thiểu hóa gradient âm của hàm mất mát trên phần dư, tạo ra dự đoán tích lũy:
\begin{equation}
    \hat{y}_i^{(t)} = \hat{y}_i^{(t-1)} + \eta \cdot f_t(\mathbf{x}_i)
\end{equation}
trong đó $\eta$ là hệ số thu hẹp (shrinkage/learning rate) kiểm soát đóng góp của mỗi cây, giúp chống overfitting \cite{Chen_2016}.

\textbf{XGBoost} (eXtreme Gradient Boosting) là một triển khai tối ưu của Gradient Boosting do Chen và Guestrin (2016) đề xuất, sử dụng cây quyết định (CART) làm weak learner. Hàm mục tiêu của XGBoost bao gồm hai thành phần:
\begin{equation}
    \mathcal{L}(\boldsymbol{\theta}) = \sum_{i=1}^{n} l(y_i, \hat{y}_i) + \sum_{t=1}^{T} \Omega(f_t)
\end{equation}
trong đó $l(y_i, \hat{y}_i)$ là hàm mất mát (thường là squared error cho bài toán hồi quy) đo sai lệch giữa giá trị thực và dự đoán, và $\Omega(f_t) = \gamma \cdot |T_{\text{leaves}}| + \frac{1}{2}\lambda \sum_{j=1}^{|T_{\text{leaves}}|} w_j^2$ là hạng tử chính quy hóa (regularization) phạt độ phức tạp của cây thông qua số lá $|T_{\text{leaves}}|$ và giá trị trọng số $w_j$ tại mỗi lá \cite{Chen_2016}. Hai siêu tham số $\gamma$ (minimum loss reduction) và $\lambda$ (L2 regularization) cho phép kiểm soát sự cân bằng giữa khả năng khớp dữ liệu và tính tổng quát.

Quá trình tìm điểm phân tách tối ưu (split finding) trong XGBoost sử dụng xấp xỉ Taylor bậc hai của hàm mục tiêu. Tại mỗi nút, thuật toán tính gradient bậc nhất $g_i = \partial l / \partial \hat{y}_i$ và gradient bậc hai $h_i = \partial^2 l / \partial \hat{y}_i^2$, sau đó chọn phân tách tối đa hóa mức giảm mất mát (gain):
\begin{equation}
    \text{Gain} = \frac{1}{2}\left[\frac{G_L^2}{H_L + \lambda} + \frac{G_R^2}{H_R + \lambda} - \frac{(G_L + G_R)^2}{H_L + H_R + \lambda}\right] - \gamma
\end{equation}
trong đó $G_L, H_L$ và $G_R, H_R$ lần lượt là tổng gradient bậc nhất và bậc hai của tập mẫu thuộc nhánh trái và phải \cite{Chen_2016}. Ngoài ra, XGBoost triển khai kỹ thuật \textbf{column subsampling} (lấy mẫu ngẫu nhiên theo đặc trưng) tương tự Random Forest, giúp giảm overfitting và tăng tốc huấn luyện.

XGBoost sở hữu nhiều ưu điểm phù hợp cho bài toán dự báo chất lượng không khí. \textit{Thứ nhất}, mô hình xử lý hiệu quả các đặc trưng không đồng nhất — bao gồm biến liên tục (nhiệt độ, tốc độ gió), biến rời rạc (giờ, ngày trong tuần) và biến có phân phối khác nhau — mà không yêu cầu chuẩn hóa đầu vào \cite{Chen_2016}. \textit{Thứ hai}, cơ chế early stopping cho phép dừng huấn luyện khi hiệu suất trên tập validation không cải thiện, tránh overfitting trên dữ liệu nhiễu vốn phổ biến trong dữ liệu quan trắc môi trường. \textit{Thứ ba}, XGBoost cung cấp chỉ số feature importance dựa trên tần suất sử dụng và mức gain trung bình, hỗ trợ diễn giải mô hình (interpretability) — một yếu tố quan trọng trong nghiên cứu môi trường \cite{tirink2025aqi}. Nghiên cứu của Tırınk (2025) trên dữ liệu AQI 8 năm tại Thổ Nhĩ Kỳ cho thấy XGBoost đạt $R^2 = 0{,}999$ và $RMSE = 0{,}234$, vượt trội so với LightGBM và SVM \cite{tirink2025aqi}. Al-Nuaimi et al. (2025) còn chứng minh XGBoost có thể nén thành mô hình ``tiny'' chỉ 13 kB với mức giảm $R^2$ khoảng 0,07, mở ra khả năng triển khai trên thiết bị IoT biên \cite{alnuaimi2025lightweight}.

Mô hình XGBoost trong hệ thống được điều khiển bởi tập siêu tham số chuyên biệt. Để đảm bảo tính chặt chẽ trong quá trình biểu diễn, các siêu tham số được quy hoạch thành hai khối cấu trúc chính:

\paragraph{Không gian Siêu tham số (Hyperparameter Space):}

\begin{table}[htbp]
\centering
\caption{Siêu tham số khối cấu trúc nền tảng của XGBoost}
\label{tab:xgb_structure}
\renewcommand{\arraystretch}{1.3}
\begin{tabular}{p{4.2cm} p{10.2cm}}
\toprule
\textbf{Siêu tham số} & \textbf{Ý nghĩa} \\
\midrule
$n\_estimators$ & Số lượng vòng lặp tăng cường tối đa, quyết định giới hạn số cây hồi quy thành phần được khởi tạo. \\
$max\_depth$ & Độ sâu tối đa phân vùng của một cây quyết định, giới hạn cấp bậc tương tác phi tuyến của các đặc trưng. \\
$min\_child\_weight$ & Tổng giá trị đạo hàm bậc hai (Hessian) tối thiểu yêu cầu tại một nút lá, nhằm tạo điều kiện dừng phân tách sớm trên các phân nhóm nhiễu. \\
\bottomrule
\end{tabular}
\end{table}

\begin{table}[htbp]
\centering
\caption{Siêu tham số khối tối ưu hóa và hệ thống của XGBoost}
\label{tab:xgb_optim}
\renewcommand{\arraystretch}{1.3}
\begin{tabular}{p{4.2cm} p{10.2cm}}
\toprule
\textbf{Siêu tham số} & \textbf{Ý nghĩa} \\
\midrule
$learning\_rate$ & Hệ số co hẹp kiểm soát gia số cập nhật của hàm mục tiêu tại từng chu kỳ tăng cường, định hình tốc độ đường phân rã hội tụ. \\
$subsample$ & Tỷ lệ trích mẫu ngẫu nhiên không hoàn lại đối với tập dữ liệu quan trắc phục vụ xây dựng cây mới, điều tiết mức phương sai nội tại. \\
$colsample\_bytree$ & Tỷ lệ lấy mẫu đặc trưng ngẫu nhiên trong mỗi phiên khởi tạo cây, hạn chế hiện tượng phụ thuộc quá hồi quy vào biến độc lập chi phối. \\
$early\_stopping\_rounds$ & Tham số xác định số vòng tăng cường liên tiếp không ghi nhận xu hướng thu hẹp của hàm tổn thất trên tập kiểm định, phục vụ cơ chế ngắt huấn luyện. \\
$eval\_metric$ & Tiêu chuẩn sai số toàn phương trung bình sử dụng làm hàm lượng giá quá trình hội tụ trong tiến trình hồi quy. \\
$reg\_alpha$ & Trọng số chính quy hóa L1 áp dụng lên hàm mục tiêu, tạo tính thưa bằng cách thu hẹp các trọng số phân nhánh không hiệu quả về bằng không. \\
$reg\_lambda$ & Trọng số chính quy hóa L2 trên lá cây, kiểm soát độ nhẵn của phép ước lượng, ngăn chặn rủi ro độ lớn tham số tiến về cực đoan. \\
$tree\_method$ & Thuật toán lưới tổ chức phân tách rời rạc hóa đặc trưng, thiết lập giới hạn tốc độ tìm kiếm nhánh. \\
$device$ & Định tuyến cấu hình phần cứng xử lý thuật toán song song hóa. \\
$n\_jobs$ & Giới hạn cấp phát luồng trên kiến trúc đa lõi trung tâm định hình số tác vụ tương tranh. \\
$random\_state$ & Hằng số khởi tạo định dạng hệ sinh số giả ngẫu nhiên đảm bảo tính lặp lại của thực nghiệm. \\
\bottomrule
\end{tabular}
\end{table}

Ngoài các siêu tham số trên, pipeline XGBoost còn bao gồm các tham số thiết kế đặc trưng  ảnh hưởng gián tiếp đến hiệu suất mô hình: cấu hình lag ($\{1, 2, 3, 6, 12, 24, 36, 48, 60, 72\}$ bước), cửa sổ rolling ($\{3, 6, 12, 24, 48\}$ bước), số lượng trạm lân cận $k$ trong thuật toán KNN không gian, và danh sách biến khí tượng tương lai được tích hợp theo từng tầm nhìn dự báo. Toàn bộ giá trị cụ thể được sử dụng trong thực nghiệm sẽ được trình bày chi tiết tại Chương 2.



\subsection{Nhóm mô hình Deep Learning — Temporal}
Trong lĩnh vực dự báo chuỗi thời gian, các mô hình Deep Learning temporal đã phát triển qua nhiều thế hệ kiến trúc. Thế hệ đầu tiên dựa trên mạng hồi quy (Recurrent Neural Networks), đại diện là LSTM với cơ chế ba cổng (quên, đầu vào, đầu ra) và GRU với cấu trúc hai cổng đơn giản hơn, cho phép học các phụ thuộc dài hạn trong chuỗi \cite{silveira2025lstmgru}. Các biến thể như BiLSTM-GRU \cite{he2024bilstmgru} và CNN-LSTM \cite{zhang2025cnnlstm} tiếp tục cải thiện hiệu suất bằng cách kết hợp đa chiều hoặc tích hợp trích xuất đặc trưng cục bộ. Tuy nhiên, nhóm mô hình RNN có hai hạn chế cố hữu: (i) xử lý tuần tự (sequential processing) khiến tốc độ huấn luyện chậm và khó song song hóa, và (ii) hiệu suất suy giảm trên chuỗi rất dài do thông tin bị ``nén'' qua bottleneck hidden state. Những hạn chế này thúc đẩy sự phát triển của các kiến trúc phi hồi quy — bao gồm MLP-based và Transformer-based — vốn xử lý toàn bộ cửa sổ đầu vào đồng thời và đạt hiệu suất cạnh tranh hoặc vượt trội trên nhiều benchmark.
\subsubsection{XLinear}
XLinear là mô hình dự báo chuỗi thời gian dài hạn dựa trên MLP (Multi-Layer Perceptron) do Chen et al. (2026) đề xuất, được thiết kế đặc biệt cho bài toán có biến ngoại sinh (exogenous inputs) \cite{Chen_2026}. Kiến trúc XLinear tổ chức luồng dữ liệu thành hai đường dẫn song song:
\begin{itemize}
    \item \textit{Đường dẫn nội sinh} (Endogenous pathway): xử lý biến mục tiêu ($PM_{2.5}$), sử dụng phép chiếu tuyến tính (linear projection) ánh xạ trực tiếp cửa sổ đầu vào $L$ sang cửa sổ dự báo $H$.
    \item \textit{Đường dẫn ngoại sinh} (Exogenous pathway): xử lý các biến khí tượng (nhiệt độ, độ ẩm, gió, áp suất), sử dụng \textbf{Gating Block} để chọn lọc thông tin hữu ích.
\end{itemize}
Gating Block là thành phần cốt lõi của XLinear, hoạt động theo cơ chế:
\begin{equation}
    \text{Gating}(\mathbf{x}) = \sigma(W_g \cdot \mathbf{x} + b_g) \odot (W_v \cdot \mathbf{x} + b_v)
\end{equation}
trong đó $\sigma$ là hàm sigmoid đóng vai trò ``cổng'' lọc, và phần tuyến tính cung cấp giá trị đặc trưng. Cơ chế này cho phép mô hình tự động đánh giá mức độ đóng góp của từng biến ngoại sinh đối với dự báo, loại bỏ nhiễu từ các biến không liên quan \cite{Chen_2026}. XLinear sử dụng hàm mất mát AsymmetricHuberLoss, phạt nặng hơn đối với sai số dự đoán thấp (under-prediction) so với dự đoán cao — phù hợp trong bài toán dự báo ô nhiễm, nơi việc bỏ sót đợt ô nhiễm nghiêm trọng nguy hiểm hơn cảnh báo sai. Ưu điểm nổi bật của XLinear là mô hình nhẹ (lightweight), tốc độ huấn luyện nhanh, và đạt hiệu suất cạnh tranh hoặc vượt trội so với các mô hình Transformer phức tạp hơn trên nhiều benchmark chuỗi thời gian \cite{Chen_2026}.

\paragraph{Không gian Siêu tham số (Hyperparameter Space):}
Mô hình XLinear cấu thành từ tập hợp các thông số định lượng được phân bổ theo hai khối chính:

\begin{table}[htbp]
\centering
\caption{Siêu tham số khối cấu trúc mạng nội tại của XLinear}
\label{tab:xlinear_structure}
\renewcommand{\arraystretch}{1.3}
\begin{tabular}{p{3.8cm} p{10.5cm}}
\toprule
\textbf{Siêu tham số} & \textbf{Ý nghĩa} \\
\midrule
$seq\_len$ & Kích thước không gian của cửa sổ thời gian lịch sử được đưa vào mô hình dưới dạng ma trận đầu vào. \\
$pred\_len$ & Tầm ngắm đích xác định giới hạn không gian thời gian của vector suy diễn kết quả đầu ra dự báo. \\
$enc\_in$ & Số lượng biến đo lường chuỗi thời gian đưa vào hệ thống mã hóa đầu cấp. \\
$features$ & Giao thức phân rã kênh luồng thời gian, kiểm soát chiến thuật ánh xạ từ ma trận không gian đa biến về đầu tính toán đơn biến đích. \\
$d\_model$ & Kích thước không gian mã hóa chỉ định lượng chiều của vector biểu diễn tuyến tính tổng cục. \\
$t\_ff$ & Kích thước không gian mạng nơ-ron lan truyền tiến đóng vai trò phân tích luật phân phối tương quan nội sinh tại mô đun thời gian. \\
$c\_ff$ & Kích thước không gian mạng lan truyền tiến ẩn quản trị phép đối chiếu liên kết tương hỗ giữa các đặc trưng dạng đa kênh tại mô đun biến ngoại sinh. \\
$embed\_dropout$ & Trọng số xác suất vô hiệu hóa đầu ra độc lập ngay sau quá trình ánh xạ không gian của biến giả lập tuyến tính ban đầu. \\
$head\_dropout$ & Trọng số xác suất thiết lập rào cản rơi cục bộ trước mạng biến đổi ánh xạ dữ liệu trực diện hướng đầu ra cuối cùng. \\
$t\_dropout$ & Trọng số xác suất ngắt kết nối thiết lập tại cụm khối van điều tiết tương quan tuần tự lịch sử thời gian. \\
$c\_dropout$ & Trọng số xác suất thiết lập thao tác chốt van tại cụm khối đối chiếu luật phân phối đặc trưng ngoại đệm. \\
$usenorm$ & Tham số logic khởi động quy trình chuẩn hóa chuỗi dữ liệu gốc thành nghịch đảo phương sai nhằm cô lập các đặc tính bất ổn định nội tĩnh toàn hệ. \\
\bottomrule
\end{tabular}
\end{table}

\begin{table}[htbp]
\centering
\caption{Siêu tham số khối huấn luyện và hàm mất mát của XLinear}
\label{tab:xlinear_optim}
\renewcommand{\arraystretch}{1.3}
\begin{tabular}{p{3.8cm} p{10.5cm}}
\toprule
\textbf{Siêu tham số} & \textbf{Ý nghĩa} \\
\midrule
$learning\_rate$ & Độ lớn bước nhảy điều tiết tốc độ trượt vi phân trên mặt cong không gian của thuật toán tối ưu gốc AdamW. \\
$weight\_decay$ & Hệ số cưỡng bức lùi định hình chuẩn L2 nhằm phá vỡ sự bùng nổ phương sai của trọng số mạng hồi quy. \\
$T\_max$ & Biên độ kích thước chuỗi nửa chặng của bộ lịch trình giảm nhiệt lượng tính theo hàm lượng giác định hướng. \\
$eta\_min$ & Giá trị cận dưới giới hạn của đường tiệm cận phân rã bước nhảy suy hao không triệt tiêu đồ hình hàm gradient. \\
$delta$ & Ngưỡng tham chiếu thiết lập điểm gãy dịch chuyển giới hạn luật cấu trúc trừng phạt chuyển trạng thái từ đa thức bậc hai thành dạng tuyến tính bất biến. \\
$alpha$ & Trọng số biểu thị tính chịu lỗi bất đối xứng quy ước việc gia tăng rủi ro tổn thất tại tình thế quan trắc chênh lệch phía dưới thấp đối chiếu giá trị dự báo thực. \\
$epochs$ & Số lượng vòng lặp cực đại được bộ giả lập đưa toàn hệ lưới dữ liệu chạy bao trọn diện tích tính toán mạng hàm. \\
$batch\_size$ & Khối lượng sức chứa dung tích nạp khối cụm tập mẫu vào trong mạch phần cứng xử lý song hành phân lớp tính. \\
$patience$ & Giới hạn khóa chốt điều kiện không ghi nhận biểu hiện hàm suy giảm hội tụ trong khoảng thời gian đủ trễ, buộc phong tỏa quy trình xử lý học để bảo tồn tiến trình hội tụ chung. \\
$corr\_threshold$ & Ngưỡng định lượng Pearson đo đếm tham số đặc trưng tối thiểu để lưới sàng lọc kích hoạt lấy đầu vào đại số hợp thức. \\
\bottomrule
\end{tabular}
\end{table}



\subsubsection{iTransformer}

iTransformer (Inverted Transformer) là kiến trúc Transformer được Liu et al. (2024) thiết kế lại cho bài toán dự báo chuỗi thời gian đa biến \cite{liu2024itransformer}. Khác với Transformer truyền thống — vốn áp dụng cơ chế attention trên chiều thời gian (mỗi bước thời gian là một token) — iTransformer thực hiện \textbf{đảo ngược} (inversion): mỗi \textit{biến đặc trưng} (feature/channel) được coi là một token độc lập. Cụ thể, cho ma trận đầu vào $\mathbf{X} \in \mathbb{R}^{L \times d}$ với $L$ bước thời gian và $d$ đặc trưng, iTransformer thực hiện:
\begin{enumerate}
    \item \textit{Embedding}: mỗi chuỗi đơn biến $\mathbf{x}^{(j)} \in \mathbb{R}^{L}$ (với $j = 1, \ldots, d$) được ánh xạ thành vector embedding thông qua phép chiếu tuyến tính, tạo thành $d$ token.
    \item \textit{Attention}: cơ chế self-attention hoạt động trên $d$ token, học mối tương quan giữa các biến (cross-variate dependency) — ví dụ: mối quan hệ giữa $PM_{2.5}$, nhiệt độ và tốc độ gió.
    \item \textit{Feed-Forward}: mạng MLP xử lý từng token để mã hóa thông tin thời gian trong mỗi biến.
    \item \textit{Projection}: ánh xạ embedding trở lại không gian dự báo $\mathbb{R}^{H}$ cho mỗi biến.
\end{enumerate}
Thiết kế đảo ngược này mang lại hai lợi thế quan trọng: (i) attention giữa các biến cho phép mô hình nắm bắt tương quan phức tạp giữa $PM_{2.5}$ và các yếu tố khí tượng mà không cần feature engineering thủ công; và (ii) MLP trên chiều thời gian xử lý hiệu quả chuỗi dài mà không gặp vấn đề độ phức tạp $O(L^2)$ của temporal attention \cite{liu2024itransformer}. Trên các benchmark chuỗi thời gian đa biến chuẩn, iTransformer đạt hiệu suất vượt trội so với các biến thể Transformer trước đó (Autoformer, FEDformer, PatchTST), đồng thời duy trì chi phí tính toán hợp lý \cite{liu2024itransformer}.

\paragraph{Không gian Siêu tham số (Hyperparameter Space):}

\begin{table}[htbp]
\centering
\caption{Siêu tham số khối cấu trúc mạng nội tại của iTransformer}
\label{tab:itrans_structure}
\renewcommand{\arraystretch}{1.3}
\begin{tabular}{p{3.8cm} p{10.5cm}}
\toprule
\textbf{Siêu tham số} & \textbf{Ý nghĩa} \\
\midrule
$seq\_len$ & Không gian ma trận cấp trước quy định lượng bản ghi lịch sử phục vụ việc tính toán phân tích ban ngang. \\
$pred\_len$ & Tọa độ ma trận tầm vươn giới hạn không gian điểm đích số liệu tương lai cần thiết lập tính dự toán. \\
$enc\_in$ & Tổng mức kênh độc lập đo quy định sự biến thiên chuyển đổi nghịch đảo ma trận tạo ra token biến định dạng. \\
$d\_model$ & Kích thước không gian cấp phát số chiều biểu diễn đặc trưng phục vụ phân tách nội suy trên một dạng token nhận tính. \\
$n\_heads$ & Số lượng hướng phân mảnh cấp không gian con thiết lập vận hành mạng thuật toán tự chú ý ma trận trực tiếp hệ Token đa hướng. \\
$e\_layers$ & Chuẩn đo số thứ bậc chồng tầng phân rã chuỗi thuật toán mã hóa Transformer dạng Inverted tuần hoàn đa hợp nối tiếp nhau. \\
$d\_ff$ & Số lượng nút trung lộ hệ mạng ẩn phi tuyến đan mạch cấu thành cấu trúc lan truyền thuận một phía nằm cục bộ giữa mỗi tháp Transformer. \\
$activation$ & Luật phân bố đạo hàm kích hoạt hàm suy diễn thiết lập ranh giới định hình thay đổi của một cụm số tuyến lưới đồ thị ngã rẽ. \\
$dropout$ & Xác suất định tuyến ngắt liên lạc ngẫu nhiên nút nhánh rèm rèn khả năng sức bền tổng duyệt chéo trên nhiều dữ liệu ẩn số mới cấu trúc tổng thể. \\
$use\_norm$ & Tiêu chuẩn logic quy định hoạt động cấu trúc RevIN áp giải các trạng thái chuổi dữ kiện dị biệt có tính độ nhiễu dao động lớn theo tỷ lệ hệ bình phân chu kỳ đảo. \\
\bottomrule
\end{tabular}
\end{table}

\begin{table}[htbp]
\centering
\caption{Siêu tham số khối huấn luyện và hàm mất mát của iTransformer}
\label{tab:itrans_optim}
\renewcommand{\arraystretch}{1.3}
\begin{tabular}{p{3.8cm} p{10.5cm}}
\toprule
\textbf{Siêu tham số} & \textbf{Ý nghĩa} \\
\midrule
$learning\_rate$ & Chỉ số trọng yếu điều khiển nhịp bước vi phân gradient quy định tốc độ thuật toán AdamW truy suất cực trị. \\
$weight\_decay$ & Thuộc tính hệ định hình biên áp phạt lùi mức độ tăng trưởng vượt trội giá trị thực số của bộ biểu mô mạng liên kết. \\
$T\_max$ & Chỉ số xác định độ mở vòng cung tính dao động giảm chấn biên giới chu kì hàm sin lịch trình hảm bước trượt vi phân tập chuẩn hóa. \\
$eta\_min$ & Trần tựa chống đóng băng giá trị bước tỷ lệ học duy trì giới hạn tối giản giữ liên tiếp hội tụ quy trình đi lại. \\
$alpha$ & Trọng lượng trung bình lai hệ thống hóa thiết lập hàm lượng đánh giá trực tiếp tiêu chí sai số tuyệt đối lẫn luật chuẩn bình phương phân kỳ cực trị trong luật đánh giá rủi ro. \\
$epochs$ & Tổng thước đo chu trình duyệt phủ tuần tự dải toàn tệp dữ liệu lưu trữ đồ thị mạng. \\
$batch\_size$ & Giới hạn luồng dung trữ mẫu xử lý chuyển pha tương tác chia cụm tại vi mạch điện toán điện trường GPU. \\
$patience$ & Thuộc tính định dạng tiêu chí khóa mạch khóa chu kì hệ hãm giám thị giới tuyến mức tiệm cận tối thiểu chốt chặt điểm cực không tối ưu mới rơi trạng thái quá cấu trúc hàm. \\
\bottomrule
\end{tabular}
\end{table}

\subsection{Nhóm mô hình Spatio-Temporal}

Các mô hình trình bày ở mục 1.2.2 và 1.2.3 khai thác phụ thuộc thời gian (temporal dependency) nhưng xử lý mỗi trạm quan trắc một cách độc lập, chưa mô hình hóa mối quan hệ không gian giữa các trạm. Trong thực tế, nồng độ $PM_{2.5}$ tại một vị trí chịu ảnh hưởng mạnh bởi phát thải và điều kiện khí tượng từ các khu vực lân cận thông qua quá trình vận chuyển khí quyển \cite{dam2024pm25}. Để khai thác phụ thuộc không gian này, mạng lưới $N$ trạm quan trắc được biểu diễn dưới dạng \textbf{đồ thị} $\mathcal{G} = (\mathcal{V}, \mathcal{E}, \mathbf{A})$, trong đó mỗi đỉnh là một trạm, mỗi cạnh biểu diễn mối quan hệ không gian, và ma trận kề $\mathbf{A} \in \mathbb{R}^{N \times N}$ mã hóa cường độ liên kết giữa các trạm \cite{panja2025estgcn}.

Trên nền tảng biểu diễn đồ thị, phép tích chập đồ thị (Graph Convolution) cho phép cập nhật biểu diễn đặc trưng của mỗi nút bằng cách \textbf{tổng hợp thông tin từ các nút lân cận} (message passing): biểu diễn $\mathbf{H}$ được biến đổi tuyến tính qua ma trận trọng số $\mathbf{W}$, sau đó nhân với ma trận kề $\mathbf{A}$ để tổng hợp thông tin không gian, tức $\mathbf{H}' = \phi(\mathbf{A} \cdot \mathbf{H} \cdot \mathbf{W})$ \cite{panja2025estgcn}. Phép tích chập đồ thị đóng vai trò thành phần nền tảng trong cả hai mô hình spatio-temporal được sử dụng trong đồ án: E-STGCN và ST-XLinear.

\subsubsection{E-STGCN}

E-STGCN (Enhanced Spatio-Temporal Graph Convolutional Network) là mô hình kết hợp phép tích chập đồ thị và LSTM để khai thác đồng thời phụ thuộc không gian lẫn thời gian trong dữ liệu đa trạm, được Panja et al. (2025) đề xuất cho bài toán dự báo chất lượng không khí \cite{panja2025estgcn}. Kiến trúc E-STGCN xử lý qua hai giai đoạn nối tiếp:

\textit{(i) Giai đoạn không gian --- GraphConv:} Tại mỗi bước thời gian $t$, ma trận đặc trưng $\mathbf{X}_t \in \mathbb{R}^{N \times F}$ (với $N$ trạm, $F$ đặc trưng) được biến đổi tuyến tính, sau đó nhân với ma trận kề $\mathbf{A}$ để tổng hợp thông tin từ các trạm lân cận. Kết quả được kết hợp với biểu diễn gốc theo chiến lược nối vector (concatenation) và kích hoạt bằng hàm ReLU:
\begin{equation}
    \mathbf{H}_t = \text{ReLU}\left(\left[\mathbf{X}_t\mathbf{W} \;\|\; \mathbf{A} \cdot \mathbf{X}_t\mathbf{W}\right]\right)
\end{equation}
trong đó $\|$ ký hiệu phép nối, tạo ra biểu diễn tích hợp cả thông tin nội tại lẫn thông tin từ lân cận không gian.

\textit{(ii) Giai đoạn thời gian --- LSTM:} Chuỗi đầu ra $\{\mathbf{H}_1, \ldots, \mathbf{H}_T\}$ của GraphConv qua $T$ bước thời gian được đưa vào mạng LSTM xử lý độc lập cho từng nút. Trạng thái ẩn cuối cùng $h_T$ tổng hợp toàn bộ thông tin lịch sử, sau đó ánh xạ qua lớp tuyến tính sang cửa sổ dự báo \cite{panja2025estgcn}.

Ma trận kề trong E-STGCN được xây dựng \textbf{một lần} trước khi huấn luyện, sử dụng Gaussian kernel trên khoảng cách Haversine $d_{ij}$ giữa các trạm:
\begin{equation}
    A_{ij} = \exp\left(-\frac{d_{ij}^2}{\sigma^2}\right)
\end{equation}
trong đó $\sigma^2$ là phương sai thực nghiệm, sau đó chuẩn hóa theo hàng. Đây là ma trận kề \textbf{tĩnh}, giả định mối quan hệ không gian cố định theo thời gian. Hạn chế chính của cách tiếp cận này là chưa phản ánh tính động của quá trình vận chuyển ô nhiễm: trong thực tế, cường độ lan truyền phụ thuộc mạnh vào hướng gió và vận tốc gió tại từng thời điểm — trạm xuôi gió ảnh hưởng lớn hơn so với trạm ngược gió, ngay cả khi ở cùng khoảng cách \cite{panja2025estgcn}.

\paragraph{Không gian Siêu tham số (Hyperparameter Space):}
Mô hình E-STGCN sử dụng một tập hợp các siêu tham số nhằm quản lý độ phức tạp thông tin giữa luồng không gian và luồng thời gian, bao gồm:

\begin{table}[htbp]
\centering
\caption{Siêu tham số khối cấu trúc mạng nội tại của E-STGCN}
\label{tab:estgcn_structure}
\renewcommand{\arraystretch}{1.3}
\begin{tabular}{p{4cm} p{10.3cm}}
\toprule
\textbf{Siêu tham số} & \textbf{Ý nghĩa} \\
\midrule
$seq\_len$ & Trục thời gian ấn định lượng mẫu cửa sổ thông tin liên luồng cấp một của hệ phân cực quá khứ. \\
$pred\_len$ & Không gian cấu trúc một điểm kết giới tầm quan sát nhịp xung đa trạm xuất chung đồng thời định lý tương lai. \\
$out\_feat$ & Trọng số đa chiều tham chiếu kích thước kênh không gian tịnh tiến sản sinh sau quá trình nhân chập hàm biểu kiến tính diện tích đồ thị mạng không gian lân cận địa đồ. \\
$lstm\_units$ & Mức định số nơ-ron trạng thái trong thành phần kiến trúc nhớ duy trì tín hiệu hồi quy tuyến không gian vận tốc thông tin diễu lịch. \\
$combination\_type$ & Thuật logic đại số định hướng phương án tích hợp và đối lưu mảng thông tin điểm quy trạm tương đối cục bộ đối lưu ma trận dữ kiện liên tuyến lưới trạm giáp ranh mở khối lượng. \\
$sigma2$ & Đại lượng hệ phương sai định chuẩn hàm định phân phối Gauss thiết lập khoảng phân rã ranh giới định hệ làm trơn lực liên kết ma trận khoảng cách tương đối quy chuẩn mặt cắt vĩ mô địa. \\
\bottomrule
\end{tabular}
\end{table}

\begin{table}[htbp]
\centering
\caption{Siêu tham số khối huấn luyện và hàm mất mát của E-STGCN}
\label{tab:estgcn_optim}
\renewcommand{\arraystretch}{1.3}
\begin{tabular}{p{4cm} p{10.3cm}}
\toprule
\textbf{Siêu tham số} & \textbf{Ý nghĩa} \\
\midrule
$learning\_rate$ & Hệ định phân rã tốc độ bước trượt thiết lập điểm dò độ hàm dốc thuật trị hướng trục AdamW cực định mặt cong phẳng điểm lõm đáy lưới tiêu chuẩn hàm học hồi quy. \\
$weight\_decay$ & Trọng số cưỡng chế tỷ lượng hàm phạt lùi mức co hẹp của hệ thống phân luồng chặn sự dư vi mạng lưới nút liên tỏa. \\
$delta$ & Đặc tính tiêu điểm phân biệt dịch chuyển trạng thái tính đạo hàm giới hạn tổn hạn phạt bình phương tuyến quy chiếu hệ L1 L2 tại trung giới điểm mất cân trượt đỉnh cao nhọn quá vi. \\
$epochs$ & Giới mức vòng chạy xuyên suốt tuyến chu kì tuần tri toàn lưới bộ định khuôn tổng hợp mẫu điểm chuẩn ranh. \\
$batch\_size$ & Tính lượng biên tải cho phép gom tụ phân cụm khối số trích mẫu chạy vòng liên khối trượt nạp phân hóa đồ hình xung thuật. \\
$patience$ & Chốt tham chiếu chặn điều kiện rơi quy trạng lặp bão hòa không truy xuất tín hiệu gia tốc tốt hạn ranh điểm chốt suy giám sai kiểm định vòng khép. \\
\bottomrule
\end{tabular}
\end{table}


\subsubsection{Mô hình ST-XLinear đề xuất}

\begin{figure}[htbp]
    \centering 
    \includegraphics[scale=0.3]{image/Kiến trúc tổng quan ST-XLinear.drawio.png} 
    \caption{Kiến trúc tổng quan của mô hình đề xuất ST-XLinear}
    \label{fig:tongquanST-XLinear}
\end{figure}

ST-XLinear (Spatio-Temporal XLinear) là mô hình mạng thần kinh nhân tạo do đồ án đề xuất, kết hợp sức mạnh phân tích thời gian của kiến trúc XLinear \cite{Chen_2026} với cơ chế học biểu diễn không gian thông qua module tích chập đồ thị động (Dynamic Graph Builder). Kiến trúc này được thiết kế nhằm khắc phục hai hạn chế lớn trong các nghiên cứu trước đây: (i) XLinear nguyên bản chỉ trích xuất đặc trưng theo chiều thời gian mà bỏ qua sự tương tác không gian giữa các trạm quan trắc, và (ii) các mô hình như E-STGCN \cite{panja2025estgcn} chỉ sử dụng ma trận kề tĩnh dựa trên khoảng cách, không phản ánh được tính động lực học của quá trình lan truyền ô nhiễm theo chiều gió.

Dựa trên Hình \ref{fig:tongquanST-XLinear}, luồng hoạt động của mô hình ST-XLinear được chia thành các giai đoạn chính như sau:

\paragraph{1. Xây dựng Ma trận kề động (Dynamic Graph Builder):} 
Thay vì sử dụng đồ thị cố định, module này liên tục tính toán ma trận kề $\mathbf{A}_t$ thay đổi theo thời gian thực dựa trên sự pha trộn giữa hai yếu tố:
\begin{itemize}
    \item \textit{Yếu tố bề mặt tĩnh (Khoảng cách)} $\mathbf{H}$: Khoảng cách vật lý (Haversine) giữa mọi cặp trạm $i$ và $j$ được chuyển hóa qua hàm Gaussian kernel: $H_{ij} = \exp(-d_{ij}^2/\sigma^2)$, thiết lập một nền tảng liên kết cơ sở cho các trạm gần nhau.
    \item \textit{Yếu tố động lực học (Gió)} $\mathbf{W}_t$: Tại mỗi bước thời gian $t$, góc phương vị (bearing) $\beta_{ij}$ từ trạm $i$ đến $j$ được đối chiếu với hướng gió quan trắc $\theta_i^{(t)}$. Hệ số vận chuyển ô nhiễm được tính bằng $W_{ij}^{(t)} = \frac{\cos(\beta_{ij} - \theta_i^{(t)}) + 1}{2}$. Giá trị này đạt cực đại khi gió thổi thẳng từ trạm $i$ sang $j$.
    \item \textit{Tổng hợp:} Ma trận kề động cuối cùng bám sát nguyên lý vật lý môi trường được pha trộn qua công thức:
    \begin{equation}
        \mathbf{A}_t = \mathbf{H} \odot \left[(1 - \alpha) + \alpha \cdot \mathbf{W}_t\right]
    \end{equation}
\end{itemize}

\paragraph{2. Trích xuất Đặc trưng Không - Thời gian:}
Chuỗi cấu trúc dữ liệu đầu vào $\mathbf{X}$ (gồm biến mục tiêu $PM_{2.5}$ và biến khí tượng) đi qua cơ chế \textbf{Temporal Embedding} (Phép chiếu tuyến tính độc lập cho từng kênh) nhằm chuyển hóa thông tin thời gian thành vector đặc trưng. Sau đó, thông qua lớp \textbf{Attention Pooling}, các kênh đặc trưng này được tóm tắt thành một vector đại diện duy nhất (Node Summary) cho mỗi trạm.

Node Summary kết hợp cùng ma trận kề động $\mathbf{A}_t$ đi vào lớp tích chập không gian \textbf{Spatial GCN} để tổng hợp thông tin từ các trạm lân cận:
\begin{equation}
    \mathbf{S}^{\text{spatial}} = \text{LayerNorm}\left(\text{GELU}\left(\text{Linear}\left([\mathbf{S} \;\|\; \mathbf{A}_t \cdot \mathbf{S}]\right) + \mathbf{S}\right)\right)
\end{equation}
Kết quả sinh ra một \textit{Spatial Global Token} chứa đựng cả lịch sử thời gian cục bộ lẫn ngữ cảnh không gian toàn cục.

\paragraph{3. Khối xương sống XLinear (Backbone):}
Khác với XLinear nguyên thủy sử dụng Global Token ngẫu nhiên có thể học được (learnable token), ST-XLinear ``bơm'' trực tiếp \textit{Spatial Global Token} vừa trích xuất vào luồng xử lý biến nội sinh. 
\begin{itemize}
    \item Khối \textbf{Time-wise Gating Module (TGM)} tiếp nhận biến nội sinh ($PM_{2.5}$) và ngữ cảnh không gian này để đánh giá tầm quan trọng của các mốc thời gian trong quá khứ.
    \item Ngữ cảnh đầu ra của TGM tiếp tục làm thông tin dẫn đường (guide context) đi vào khối \textbf{Variable-wise Gating Module (VGM)} để tiến hành sàng lọc nhiễu, quyết định biến ngoại sinh (nhiệt độ, độ ẩm, v.v.) nào có ý nghĩa lớn nhất đối với dự báo.
\end{itemize}
Hai đầu ra cuối cùng được nối lại (Concatenation) và đi qua lớp \textbf{Prediction Head} để đưa ra kết quả nồng độ $PM_{2.5}$ tại từng trạm.

Dưới đây là phần mô tả siêu tham số đã được hệ thống hóa, tập trung vào ý nghĩa và vai trò của từng siêu tham số đối với quá trình định lượng mô hình:

\paragraph{Không gian Siêu tham số (Hyperparameter Space):}

\begin{table}[htbp]
\centering
\caption{Siêu tham số khối cấu trúc đồ thị mạng nội tại của ST-XLinear}
\label{tab:stxlinear_structure}
\renewcommand{\arraystretch}{1.3}
\begin{tabular}{p{3.8cm} p{10.5cm}}
\toprule
\textbf{Siêu tham số} & \textbf{Ý nghĩa} \\
\midrule
$seq\_len$ & Kích thước không gian ấn định quy chuẩn số bước thời gian thuộc mảng cấu trúc lịch sử đưa vào quá trình phân giải đại số đa chiều. \\
$pred\_len$ & Tầm ngắm cố định mục tiêu dự báo độc lập theo kiến trúc đa phiên bản thiết kế đơn đích từng độ phóng phóng điểm mốc riêng. \\
$d\_model$ & Kích thước không gian mã hóa chỉ định diện tích bề rộng cấp vector quy định ma trận biểu kiến của mảng token trạng thái tĩnh thời cực hình hệ kiến. \\
$t\_ff$ & Diện tích hệ chiều mạng hàm ẩn kết cấu tại hệ kiểm soát luồng biến thời nội sinh đánh giá sự quy định phân độ quy luật mạng nội tính toán quan trạm. \\
$c\_ff$ & Diện tích hệ lưới mạng lớp ẩn trực tuyến chuyên quản phân rã xử lý giao kênh đo quy trạm đo ngoại biến đệm khí tượng vi lân chéo kênh mảng. \\
$dropout$ & Xác suất tạo ngắt tín hiệu liên kết phá hủy quy luồng bộ khung liên tưởng hàm dẻo giảm sự học tuyến tính nguyên sinh cấu trúc phân lập cũ. \\
$alpha\_graph$ & Hệ số trọng lượng lai cố định dung hòa biên độ luân chuyển khối ma trận khoảng cách không gian thuần vật lý kết hợp đồ hình tương tác hướng gió véc-tơ đa hướng. \\
\bottomrule
\end{tabular}
\end{table}

\begin{table}[htbp]
\centering
\caption{Siêu tham số khối huấn luyện và hàm mất mát của ST-XLinear}
\label{tab:stxlinear_optim}
\renewcommand{\arraystretch}{1.3}
\begin{tabular}{p{3.8cm} p{10.5cm}}
\toprule
\textbf{Siêu tham số} & \textbf{Ý nghĩa} \\
\midrule
$learning\_rate$ & Chỉ lượng xung tạo vận dốc phân rã mức quy luồng cập nhật ma trận mạng lưới AdamW đi trục hạ hồi biên sai lệch giá trị chuẩn tập xuất. \\
$weight\_decay$ & Hệ số cưỡng đoạt giá trị triệt biên lùi trọng lượng thu hẹp quy mô dải đại số kết tỏa tuyến phân tầng liên nút khối nhân. \\
$T\_max$ & Khẩu giới tính cực chu vi đường lượn hàm cosin thiết kế hệ đếm giọt nhỏ dốc rã điểm bước tiến. \\
$eta\_min$ & Đáy đệm cấm sự biến dạng trị điểm trần hàm số hệ quy trượt giảm chặn sự lùi cực đại cực không giới vi không tồn tịnh định. \\
$mse\_w$ & Chỉ số phần tỷ định ước hàm chia pha trọng lượng tham gia hàm mức vi trừng phạt phân rẽ độ lớn tổn sai. \\
$delta$ & Vị điểm gãy hàm xác quyết lập tuyến đồ chuyển phương thức đánh giá điểm uốn vi phân của sai mức tuyến chuẩn tính ranh. \\
$epochs$ & Giới hạn mức tập lặp lưới tổng vi mô bao bọc hệ số vòng chu kì tuần hồi hàm chạy khuyết dữ liệt tập số chéo trượt gốc. \\
$batch\_size$ & Phân chia mức hộp sức dung gom trượt chứa mảng tệp liên chập vào lưới trung đới lõi nhân diện GPU tại hệ nhịp tính điện áp logic xung điện. \\
$patience$ & Định biên hãm cảnh giác báo sớm đường tiệm cận khóa hiện trạng lặp dư thừa khóa luồng đóng máy móc trói trạng thái bão hoà tĩnh định hồi quy tập thẩm âm định biên lặp suy suy biến vi ranh. \\
\bottomrule
\end{tabular}
\end{table}

\subsubsection{Phương pháp Ensemble}

Phương pháp học máy kết hợp (Ensemble Learning) là một kỹ thuật mạnh mẽ nhằm nâng cao hiệu suất dự báo và tính ổn định của hệ thống thông qua việc kết hợp dự đoán từ nhiều mô hình thành phần khác nhau. Nguyên lý cốt lõi của Ensemble dựa trên thực tế là các mô hình đơn lẻ (base learners), do sự khác biệt về kiến trúc và cách trích xuất đặc trưng, thường mắc phải các loại sai số phân bố khác nhau. Việc kết hợp chúng giúp triệt tiêu phương sai (variance), khắc phục điểm yếu cục bộ của từng mô hình và giảm thiểu rủi ro quá khớp (overfitting).

Trong hệ thống dự báo này, phương pháp Ensemble được thiết kế để kết hợp trực tiếp đầu ra của hai kiến trúc tiêu biểu có đặc thù phân tích khác nhau: mô hình dựa trên MLP (như XLinear, chuyên biệt về nắm bắt tính chu kỳ và lọc nhiễu ngoại sinh) và mô hình dựa trên đồ thị (như ESTGCN, chuyên biệt về khai thác dòng chảy không gian). Quy tắc kết hợp được sử dụng là Trung bình có trọng số (Weighted Average). Cụ thể, giá trị dự báo cuối cùng $\hat{y}$ được tính bằng chập tuyến tính giữa giá trị dự báo của hai mô hình thành phần:
\begin{equation}
    \hat{y} = \alpha \cdot \hat{y}_{XLinear} + (1-\alpha) \cdot \hat{y}_{ESTGCN}
\end{equation}
trong đó, $\hat{y}_{XLinear}$ và $\hat{y}_{ESTGCN}$ lần lượt là dự báo nồng độ $PM_{2.5}$ xuất ra từ hai mô hình, và $\alpha \in [0, 1]$ là trọng số kết hợp.

Với phương pháp luận này, $\alpha$ không phải là tham số có thể học được (learnable parameter) bằng Gradient Descent trong quá trình huấn luyện mô hình, mà được xem như một siêu tham số (hyperparameter). Để đảm bảo tính khách quan và đạt hiệu suất tối ưu, $\alpha$ được xác định thông qua \textbf{tìm kiếm dạng lưới (Grid Search)} trên tập dữ liệu kiểm chứng (Validation set). Phương pháp dò tìm này được kiểm soát bởi các tham số không gian cấu hình chuyên biệt như sau:
\begin{table}[htbp]
\centering
\caption{Siêu tham số cấu hình Grid Search của phương pháp Ensemble}
\label{tab:ensemble_params}
\renewcommand{\arraystretch}{1.3}
\begin{tabular}{p{3.8cm} p{10.5cm}}
\toprule
\textbf{Siêu tham số} & \textbf{Ý nghĩa} \\
\midrule
$alpha$ ($\alpha$) & Đại lượng vô hướng định hình tỷ lệ đóng góp của mô hình tuyến tính ST-XLinear cấu hình nội sinh trong phép toán trung bình có trọng số ghép nối luồng với hệ mạng tích chập đồ thị E-STGCN mở rộng. \\
$metric$ & Phương pháp đánh giá tiêu chuẩn thống nhất MAE thiết lập hệ sai số tiêu quan khách làm nền tảng đánh đổi kết quả truy vấn. \\
$step$ & Kích thước gia tốc tịnh tiến trong kỹ thuật tìm kiếm lưới trực giao ấn định quy mô độ mịn của thuật toán nội suy điểm quy tụ tối ưu tham số alpha hệ phân lai. \\
\bottomrule
\end{tabular}
\end{table}
Cách tiếp cận này tận dụng được điểm mạnh tối đa từ các mô hình thành phần mà không lam tăng đáng kể tài nguyên tính toán.

\subsection{Các chỉ số đánh giá hiệu suất}

Để đánh giá một cách toàn diện, khách quan độ chính xác và khả năng tổng quát hóa của các mô hình dự báo $PM_{2.5}$, bước đánh giá hiệu suất (Evaluation) đóng vai trò tối quan trọng. Không có một chỉ số đơn lẻ nào phản ánh hoàn hảo mọi khía cạnh của sai số dự báo; do đó, thông lệ nghiên cứu khoa học khuyến nghị sử dụng kết hợp nhiều chỉ số thống kê khác nhau để có cái nhìn đa chiều \cite{chai2014rmse, armstrong1992mape}. Trong đồ án này, các mô hình được so sánh chéo dựa trên bốn chỉ số đo lường chính.

\subsubsection{Root Mean Square Error (RMSE)}

Root Mean Square Error (Sai số toàn phương trung bình căn) là một trong những chỉ số đánh giá độ phân tán định lượng phổ biến nhất trong các bài toán hồi quy và dự báo môi trường. Nó được định nghĩa là căn bậc hai của trung bình cộng bình phương các khoảng cách giữa giá trị dự báo và giá trị quan trắc thực tế, có cùng đơn vị với biến mục tiêu ($\mu g/m^3$):
\begin{equation}
    RMSE = \sqrt{\frac{1}{N} \sum_{i=1}^{N} \left( y_i - \hat{y}_i \right)^2}
\end{equation}
trong đó $y_i$ là giá trị nồng độ $PM_{2.5}$ thực tế ở quan trắc thứ $i$, $\hat{y}_i$ là giá trị dự báo tương ứng, và $N$ là tổng số cỡ mẫu đánh giá.

Đặc tính quan trọng nhất của RMSE đến từ thao tác bình phương sai số trước khi lấy trung bình. Điều này khiến RMSE \textbf{phạt rất nặng (penalize)} đối với các sai số lớn (outliers) sinh ra từ việc mô hình dự đoán trật xa thực tế. Trong bối cảnh dự báo chất lượng không khí, đặc tính này cực kỳ hữu ích: việc mô hình trượt một sự kiện nồng độ $PM_{2.5}$ tăng vọt đột ngột (cực đoan) sẽ lập tức phản ánh bằng mức RMSE vọt lên cao. Theo Chai và Draxler (2014), việc sử dụng RMSE là hoàn toàn ưu việt và đáng tin cậy hơn so với các chỉ số tuyệt đối khi phân phối của sai số dự đoán tiến gần đến phân phối chuẩn Gaussian, đồng thời RMSE thỏa mãn đầy đủ bất đẳng thức tam giác để được xem là một khoảng cách (distance metric) đo lường không gian hợp lệ trong tối ưu hoá \cite{chai2014rmse}. Việc hạn chế cực tiểu hóa RMSE đồng nghĩa với việc đưa ra một cảnh báo sức khỏe an toàn hơn.

\subsubsection{Mean Absolute Error (MAE)}

Mean Absolute Error (Sai số tuyệt đối trung bình) là một chỉ số đo lường độ lớn trung bình của sai số trong tập dự báo, bỏ qua dấu của chúng. Khác với RMSE, MAE không sử dụng thao tác bình phương mà trực tiếp lấy giá trị tuyệt đối của hiệu số giữa kết quả dự báo và quan trắc thực tế, sau đó tính trung bình cộng:
\begin{equation}
    MAE = \frac{1}{N} \sum_{i=1}^{N} \left| y_i - \hat{y}_i \right|
\end{equation}

Với cấu trúc tuyến tính này, mỗi sai số đóng góp vào giá trị trung bình với một tỷ lệ thuận chính xác theo độ lớn tuyệt đối của riêng nó. Willmott và Matsuura (2005) lập luận rằng MAE là một chỉ số "tự nhiên" và ít mơ hồ (unambiguous) hơn so với RMSE \cite{willmott2005mae}. Do không cực đại hóa (phạt nặng) các điểm ngoại lai, MAE mang tính ổn định (robust) hơn trong các bộ dữ liệu có nhiều nhiễu bất thường, giúp đánh giá một cách trung lập \textbf{hiệu suất mô hình tính trên trung bình (average model performance)}. Trong khi RMSE cho biết sai số ở kịch bản xấu (tệ nhất), thì độ chênh lệch giữa MAE và RMSE (RMSE $\geq$ MAE) giúp cung cấp một cái nhìn tổng thể về phương sai của lượng sai lệch do mô hình tạo ra.

\subsubsection{Coefficient of Determination ($R^2$)}

Coefficient of Determination (Hệ số quyết định), thường được ký hiệu là $R^2$, là một chỉ số thống kê phản ánh tỷ lệ phần trăm sự biến thiên (phương sai) của biến mục tiêu mà mô hình dự báo có thể nắm bắt và giải thích được bằng các biến độc lập. Định nghĩa tổng quát của hệ số $R^2$ áp dụng cho cả hồi quy tuyến tính và phi tuyến tính (như mạng nơ-ron hay cây quyết định xgboost) \cite{nagelkerke1991r2, tirink2025aqi} được biểu diễn qua tỷ số giữa tổng bình phương phần dư ($SS_{res}$) và tổng bình phương toàn phần ($SS_{tot}$):
\begin{equation}
    R^2 = 1 - \frac{SS_{res}}{SS_{tot}} = 1 - \frac{\sum_{i=1}^{N} (y_i - \hat{y}_i)^2}{\sum_{i=1}^{N} (y_i - \bar{y})^2}
\end{equation}
trong đó, $\bar{y}$ là giá trị trung bình lượng $PM_{2.5}$ của dữ liệu thực tế đánh giá ($y_i$). 

Hệ số $R^2$ thường nhận giá trị trong khoảng $[0, 1]$. Một giá trị $R^2$ tiến gần đến $1$ chứng tỏ mô hình nội suy/ngoại suy gần như hoàn hảo, bám sát xát mọi dao động thăng trầm của dữ liệu mục tiêu. Nếu $R^2 = 0$, mô hình chỉ đơn thuần dự đoán bằng cách xuất ra giá trị trung bình $\bar{y}$ (không học được bất kỳ tri thức biểu diễn quy luật nào). Ở một số hình thức sai lệch quá nghiêm trọng từ mô hình, $R^2$ có thể âm, ám chỉ việc dự báo bằng trị số trung bình còn hiệu quả hơn chính mô hình kiến trúc đưa ra.

\subsubsection{Mean Absolute Percentage Error (MAPE)}

Mean Absolute Percentage Error (Sai số phần trăm tuyệt đối trung bình) đo lường độ chính xác dưới dạng tỷ lệ phần trăm, mang lại ưu điểm vượt trội về khả năng diễn giải trực quan cho người trực tiếp sử dụng hệ thống so với các chỉ số số tuyệt đối (metric đơn vị) như RMSE hay MAE. Công thức MAPE là bình quân của sai số tuyệt đối quy đổi về tỷ lệ phần trăm dựa trên giá trị gốc:
\begin{equation}
    MAPE = \frac{100}{N} \sum_{i=1}^{N} \left| \frac{y_i - \hat{y}_i}{y_i} \right|
\end{equation}

Mặc dù rất phổ biến và dễ sử dụng để đánh giá cũng như so sánh hiệu năng chéo giữa các loạt độ chia dataset khác nhau \cite{armstrong1992mape}, MAPE vẫn tiềm ẩn một số nhược điểm rất lớn về mặt toán học đã được Armstrong và Collopy (1992) lưu ý. \textit{Thứ nhất}, khi chuỗi thời gian chứa giá trị thực tế quan trắc ($y_i$) rất nhỏ hoặc bằng không (hoàn toàn khả dĩ trong dữ liệu nồng độ chất lượng thấp), mẫu số sẽ khiến chỉ số bị phân kỳ tiệm cận vô cực, do đó không định nghĩa được. \textit{Thứ hai}, do đặc tính công thức, MAPE có thiên hướng phạt nặng sai số phóng đại (dự đoán trượt lên cao - overpredict) mạnh hơn nhiều so với dự đoán thiếu (dự báo trượt xuống thấp - underpredict) ngay cả trên một khoảng chênh lệch tuyệt đối như nhau. Chính vì thế, trong khuôn khổ đồ án, MAPE được viện dẫn đồng thời cùng cụm metrics như MAE, RMSE để có cái nhìn tham chiếu đầy đủ.

\subsection{Các công trình nghiên cứu liên quan}

Để xây dựng cơ sở khoa học và định vị được những đóng góp của đồ án, việc rà soát và đánh giá các nghiên cứu đi trước về dự báo nồng độ ô nhiễm $PM_{2.5}$ cả trong nước và quốc tế là hết sức cần thiết. Các hướng tiếp cận đang phân hóa rất đa dạng dựa trên cả mô hình định lượng hóa học (Hóa khí quyển) truyền thống và hệ thống mạng nơ-ron học sâu (Deep Learning) hiện đại.

\subsubsection{Nghiên cứu trong nước}

Tại Việt Nam, các nghiên cứu chuyên sâu về phân bố ô nhiễm không khí và ảnh hưởng sức khỏe đang ngày càng được chú trọng do mức độ báo động tại các đô thị lớn. Khởi điểm từ khía cạnh vĩ mô, nghiên cứu của Bui và Nghiem (2025) đã chứng minh được mối liên hệ tuyến tính giữa mật độ đô thị hóa và mức độ ô nhiễm thông qua mô phỏng hồi quy trên 63 tỉnh thành: mỗi 1\% gia tăng đô thị hóa sẽ kéo theo nồng độ $PM_{2.5}$ tăng \textbf{0,54 $\mu g/m^3$} \cite{bui2025urbanization}. Cùng đánh giá trên bài toán y tế công cộng, Le et al. (2024) \cite{le2024environmental} và Hoang et al. (2023) \cite{hoang2023health} thiết lập được sự rủi ro cực đoan của bệnh lý tim mạch, hệ hô hấp với các loại bụi phát thải thứ cấp. Từ đó ước tính được quy mô kinh tế thiệt hại, và cứu sống xấp xỉ vài nghìn quy đổi sinh mạng mỗi năm nếu đạt tiêu chuẩn quản lý rủi ro như thời kỳ phong toả Covid-19. 

Về mặt kiểm soát nồng độ và khí tượng mô phỏng số lượng, thay vì trí tuệ nhân tạo, đa phần các học giả Việt Nam ưa chuộng sử dụng nền tảng mô hình vật lý - hóa học truyền thống tiêu biểu như hệ thống CMAQ (Community Multiscale Air Quality). 
Cụ thể, Ho et al. (2025) sử dụng mô hình này kết hợp với hồ sơ phát thải giao thông nhằm phân tích cơ chế lan truyền nồng độ các chất tại Hà Nội định chuẩn theo các thời vụ \cite{ho2025assessment}. Tương tự, Nguyen et al. (2024) thực hành giả lập trên khung CMAQ liên kết hệ WRF, dự đoán được biến thiên bụi độ bao phủ không gian lớn tuy nhiên đạt tỷ lệ sai số cao (thường vướng $MAPE \approx 63\%$) ở điều kiện phức tạp \cite{nguyen2024cmaq}. 
Bên cạnh đó, các nghiên cứu đi sâu vào quan trắc thẳng thực trạng vận động ô nhiễm như của Dam et al. (2024) đã hé lộ xu hướng gradient bụi lắng nghịch nhiệt theo độ cao phân tầng rõ nét: nồng độ ở 30m là 50--84 $\mu g/m^3$, trong khi ở 320m chỉ còn 34--40 $\mu g/m^3$ \cite{dam2024pm25}. 

\subsubsection{Nghiên cứu ngoài nước}

Trái ngược với sự thống trị của các mô hình cơ học hóa tại Việt Nam, xu thế nghiên cứu toàn cầu trong vài năm trở lại đây tập trung cao độ vào nhóm hệ thống máy học (Machine Learning) và mạng nơ-ron học sâu (Deep Learning).

\textbf{Ứng dụng các thuật toán kết hợp tăng cường (Boosting):} XGBoost và LightGBM thường được triển khai nhờ tốc độ xử lý nhanh, minh bạch dữ liệu đầu vào. Tırınk (2025) từng so sánh sức ép dự báo AQI trên phổ dữ liệu 8 năm tại Türkiye và kết luận XGBoost cho cấu hình hiệu năng tốt nhất, bỏ xa SVM truyền thống \cite{tirink2025aqi}. Với yêu cầu IoT ở trạm giám sát rẻ tiền biên, Al-Nuaimi et al. (2025) có thể tóm gọn một mô hình XGBoost cho cấu hình tài nguyên rất nhỏ (13 kB) mà vẫn đảm bảo được hiệu năng khả dĩ \cite{alnuaimi2025lightweight}.

\textbf{Ứng dụng mô hình học sâu cho tính chuỗi thời gian (Temporal):} Cơ chế tự hồi quy trong chuỗi được xử lý nhờ bộ ba cổng LSTM / GRU chuẩn. Theo như thiết kế lại từ Silveira et al. (2025) cho phương pháp single-step lẫn multi-horizon, kiến trúc LSTM tỏ ra vô cùng cứng cáp do năng lượng chặn đứng vấn đề triệt tiêu đạo hàm trong truyền lùi \cite{silveira2025lstmgru}. Zhang et al. (2025) \cite{zhang2025cnnlstm} đề xuất ghép CNN với LSTM để lọc đặc trưng đa kênh cục bộ ứng dụng tại Bắc Kinh (đạt $RMSE = 5{,}236$ trong biến thiên 6 giờ). Gần đây, hệ quy chiếu He et al. (2024) kết hợp LSTM hai chiều và GRU để chặn những nhiễu loạn xa gần \cite{he2024bilstmgru}. Tuy vậy, điểm nghẽn các mô hình họ RNN nói chung vướng phải vẫn là sự đền bù khổng lồ trong tính lặp lại tuyến tính mà không thể xử lý đa bước song song.

\textbf{Chuyển dịch sang mô hình không - thời gian và kiến trúc Transformer:} Để bù đắp nhược điểm về thời gian quá dài, Transformer được cải tiến (ví dụ: iTransformer \cite{liu2024itransformer}) áp dụng phương án self-attention đảo ngược kênh (channel-wise) thay thế vòng tuần hoàn của RNN. Đặc biệt với yếu tố đồ thị vận chuyển ô nhiễm, thuật toán E-STGCN (Panja et al., 2025) dùng mô hình Graph Convolution Network áp mã ma trận kề tính toán từ khoảng cách Haversine chạy Gaussian kết hợp chuỗi LSTM sau đó xử lý đồng loạt sự tập trung lan truyền $PM_{2.5}$ qua các điểm tĩnh lân cận cho những dịp cực đoan \cite{panja2025estgcn}. 

\subsubsection{Nhận xét và khoảng trống nghiên cứu}

Việc tổng hợp các tài liệu của học giả trong nước lẫn quốc tế trên nhiều phương diện (sức khỏe cộng đồng \cite{hoang2023health}, hóa môi trường \cite{ho2025assessment}, mạng học sâu nâng cao \cite{liu2024itransformer}) đã đem đến một phác hoạ đầy đủ và cũng làm bộc lộ một số góc khuất tồn đọng trong nền tảng nghiên cứu ứng dụng cho bài toán tại Việt Nam:

\begin{enumerate}
    \item \textbf{Thiếu hụt hệ thống đánh giá đối chuẩn (Benchmark) đồng thời:} Trên cùng dữ liệu quan trắc $PM_{2.5}$ tại nước ta, đa số các công trình hoặc dùng các bộ quan sát quá cũ (thập kỷ trước) để thử thuật toán Machine Learning nhỏ lẻ, hoặc dùng khung mô phỏng truyền thống tốn kém máy chủ CMAQ. Chưa có bất kỳ nghiên cứu nào tạo dựng một cuộc thi đối chuẩn trực tiếp giữa Machine Learning cổ điển, Sequence Model, Transformer và Spatio-Temporal Graph trên bộ dữ liệu Việt Nam nhiều năm với hệ quy chiếu dự báo đa khung nhìn (Multi-horizon).
    
    \item \textbf{Cấu trúc không gian đồ thị chưa khai thác đặc tính động lực học lưu chất:} Hầu hết các kỹ thuật Graph Network truyền thống trong dự báo dòng chảy ô nhiễm không khí (kể cả E-STGCN mới nhất) đều áp dụng mô hình liên kết ma trận đồ thị dạng lưới bằng khoảng cách tĩnh \cite{panja2025estgcn}. Điều này bỏ ngỏ ảnh hưởng cốt lõi của lực mang đi là "gió" đối với hạt vi nhựa sinh học và kim loại lơ lửng. Một trạm ngay sát nhưng nằm ngược dòng gió sẽ khó có thể gây lây lan ô nhiễm như một trạm ở xa tít nhưng lại trực diện theo cùng chiều gió.
    
    \item \textbf{Quá trình kỹ nghệ hóa đặc trưng (Feature Engineering) chưa dựa trên nguyên lý khí quyển học:} Điển hình do giới hạn về mặt dữ liệu, công tác thiết kế đặc trưng bổ trợ thường được bỏ qua đối với vùng nhiệt đới phức tạp tự nhiên. Chưa có thực nghiệm nào bóc tách trích xuất các dạng chỉ báo liên quan đến aerosol sulfate hay chỉ số quy luật bức xạ của ánh sáng làm lớp cung cấp bối cảnh (context guide) cho các kênh DL model Việt Nam.
\end{enumerate}

Những thiết sót nêu trên chính là cơ sở cho các đề xuất đóng góp chuyên môn học thuật cốt lõi trong toàn bộ hệ thống ST-XLinear cấu hình Dynamic Graph và kho kỹ nghệ dữ liệu mà Đồ án đang theo đuổi ở các chương tiếp theo.

\subsection{Sơ đồ tổng quan giải pháp}

Hình~\ref{fig:workflow_tongquan} trình bày sơ đồ tổng quan giải pháp của hệ thống dự báo nồng độ $PM_{2.5}$, được tổ chức thành năm khối chức năng liên kết tuần tự từ khâu thu thập dữ liệu đến khâu triển khai ứng dụng.

\begin{figure}[htbp]
    \centering 
    \includegraphics[width=\textwidth]{image/Tổng quan pipline xử lý.drawio.png} 
    \caption{Sơ đồ tổng quan giải pháp hệ thống dự báo nồng độ $PM_{2.5}$}
    \label{fig:workflow_tongquan}
\end{figure}

\textbf{Khối 1 --- Thu thập dữ liệu (Crawl Data).} Dữ liệu thô được thu thập tự động từ hai nguồn API: nền tảng Weatherbit cung cấp dữ liệu chất lượng không khí ($PM_{2.5}$, $PM_{10}$, $CO$, $O_3$, $NO_2$, $SO_2$) và nền tảng OpenMeteo cung cấp dữ liệu khí tượng lịch sử (nhiệt độ, độ ẩm, nhiệt độ điểm sương, lượng mưa, mây, tốc độ gió, hướng gió, gió giật, nhiệt độ và độ ẩm đất). Dữ liệu từ 12 trạm quan trắc trong giai đoạn 2023--2025 được lưu trữ dưới dạng tệp thô trước khi đưa vào quy trình xử lý.

\textbf{Khối 2 --- Tiền xử lý dữ liệu (Process).} Dữ liệu thô trải qua bốn giai đoạn xử lý nối tiếp, mỗi giai đoạn tạo ra sản phẩm trung gian phục vụ giai đoạn tiếp theo:

\begin{itemize}
    \item \textit{Data Cleaning:} Thực hiện ghép nối dữ liệu không khí và khí tượng, phát hiện và gắn cờ các bản ghi bất thường (dữ liệu đóng băng, lỗi cảm biến $PM_{2.5}$, giá trị ngoại lai), loại bỏ bản ghi trùng lặp, nội suy giá trị khuyết thiếu, và tạo các đặc trưng phái sinh. Sản phẩm đầu ra là các tệp dữ liệu sạch (files clean).
    
    \item \textit{Feature Engineering:} Tinh chiết các đặc trưng mới dựa trên nguyên lý hóa khí quyển (chỉ số oxidation potential, humid sulfate risk, thermal stability), phân tách thành phần gió ($\sin$/$\cos$), và trích xuất đặc trưng trượt (lag features, rolling mean, rolling std) cùng mã hóa chu kỳ thời gian.
    
    \item \textit{Normalization:} Áp dụng các chiến lược chuẩn hóa khác nhau tùy theo đặc tính phân phối của từng nhóm biến, bao gồm các tổ hợp $Log1p$, $RobustScaler$, $StandardScaler$ và $MinMaxScaler$. Các bộ biến đổi (scalers) được khớp (fit) trên tập huấn luyện và lưu trữ dưới dạng tệp .pkl để phục vụ quá trình giải chuẩn hóa khi suy luận.
    
    \item \textit{Block Split:} Phân chia dữ liệu theo cấu trúc khối thời gian tuần hoàn (cyclic block splitting) với ba cấu hình: Block~5 (chu kỳ 5 ngày: 3 ngày train, 1 ngày val, 1 ngày test), Block~7 (chu kỳ 7 ngày: 5 ngày train, 1 ngày val, 1 ngày test) và Block~30 (chu kỳ 30 ngày: 20 ngày train, 5 ngày val, 5 ngày test). Chiến lược chia theo khối tuần hoàn đảm bảo mỗi tập con đều bao phủ được các điều kiện khí tượng đa dạng qua các mùa, thay vì tập trung vào một giai đoạn thời gian cố định.
\end{itemize}

\textbf{Khối 3 --- Huấn luyện mô hình (Model Train).} Dữ liệu sau phân chia được đưa vào hệ thống đối chuẩn (benchmark) song song giữa sáu kiến trúc mô hình: XGBoost, XLinear, iTransformer, E-STGCN, ST-XLinear (mô hình đề xuất) và Ensemble. Trọng số của các mô hình sau huấn luyện được lưu trữ vào cơ sở dữ liệu mô hình (Models) để phục vụ cả quá trình đánh giá lẫn triển khai.

\textbf{Khối 4 --- Đánh giá hiệu suất (Evaluation).} Kết quả dự báo được đánh giá trên năm tầm nhìn ($T+1$, $T+3$, $T+6$, $T+12$, $T+24$ giờ), sử dụng bốn chỉ số: $RMSE$ \cite{chai2014rmse}, $MAE$ \cite{willmott2005mae}, $MAPE$ \cite{armstrong1992mape} và $R^2$ \cite{nagelkerke1991r2}. Việc đánh giá đa tầm nhìn và đa chỉ số cho phép so sánh hiệu suất của từng kiến trúc trên các kịch bản dự báo ngắn hạn đến dài hạn.

\textbf{Khối 5 --- Triển khai ứng dụng (Deployment).} Mô hình đạt hiệu suất phù hợp được nạp từ cơ sở dữ liệu mô hình vào module suy luận (Model Prediction), sau đó chuyển tiếp kết quả dự báo lên giao diện Web UI xây dựng trên nền tảng Streamlit, cung cấp khả năng giám sát chất lượng không khí và cảnh báo sớm.

Sơ đồ giải pháp trên là khung tham chiếu xuyên suốt cho các chương tiếp theo: \textbf{Chương 2} trình bày chi tiết phương tiện và phương pháp nghiên cứu tương ứng với Khối 1 đến Khối 3, trong khi \textbf{Chương 3} phân tích kết quả thực nghiệm từ Khối 4 và thảo luận về hiệu quả triển khai tại Khối 5.
